\chapter{比赛相关}
本册是当前模板所用标准库与环境内容的第一版整理，不宣称覆盖整个 C++ 标准库。尚需按实际使用记录补齐未登记的标准库细项、环境版本差异与更多经过执行的示例。

\section{提交前的接口核对}
\index{提交检查}
先对照题面确认下标、闭区间或半开区间、空输入、重边、自环和数值范围。返回值为0可能是合法答案，不能一律当作失败。算法模板的数学前提与输出方案约定优先于类型名的直觉。

本库 C++20 模板通常需要 bits/stdc++.h、cassert 和 std 命名空间。GNU 扩展不是标准 C++；若评测环境没有 libstdc++，不要假定 PBDS、rope、bits/stdc++.h 可用。整数乘积先转宽类型再乘，不能乘完才强转。断言用于检查调用前提，不能把断言本身当作正确性证明。

\section{对拍与性能检查}
\index{对拍}
小数据参考应尽量用枚举、直接模拟或不同数学公式。固定随机种子并保存失败输入；先检查样例、空边界、极值、重复值，再测试最大规模。消毒器检查越界与未定义行为，优化版单独测时间。相同机器上的本地耗时也不能直接换算为 OJ 排名。

递归 DFS 保持递归。本机默认线程栈较小时调整本机测试环境，不因此改成显式模拟 DFS。算法的栈深度与题目内存约束仍需匹配。

\chapter{工具软件}
\section{编译与运行}
\index{g++}\index{Sanitizer}
以下在 Bash 中运行；源文件和输出文件变量加引号。macOS 的 g++ 可能实际调用 Clang，先用 --version 核对；本项目在本机选择 g++-16，其他环境可设置 CXX。调试与提交使用相同语言标准，GNU 扩展用 GNU 工具链实测。
\begin{lstlisting}[language=bash]
src="main.cpp"
out="main"
g++ -std=c++20 -O2 -Wall -Wextra "$src" -o "$out"
./main < input.txt > actual.txt
diff -u expected.txt actual.txt
\end{lstlisting}
调试构建：
\begin{lstlisting}[language=bash]
g++ -std=c++20 -O1 -g -fsanitize=address,undefined \
    -fno-omit-frame-pointer main.cpp -o debug
./debug < input.txt
\end{lstlisting}
Linux 本地测试可用 \texttt{ulimit -s unlimited} 放宽允许范围内的主线程栈；macOS 使用 GNU C++ 时本库深递归测试采用链接选项 \texttt{-Wl,-stack\_size,0x20000000}。这些是本地测试配置，不应盲目加进 OJ 源文件。

\section{Bash 文件与退出状态}
\index{Bash}
\texttt{<} 重定向输入，\texttt{>} 覆盖输出文件，\texttt{>>} 追加。不要把输入文件名同时作为输出文件名。脚本中可用 \texttt{set -euo pipefail} 尽早发现命令失败、未定义变量和管道失败。\texttt{diff} 返回1通常表示内容不同，不是程序崩溃。

含空格的路径使用双引号；想保留美元符号、反引号等字面内容时，用单引号或带引号的 here-document 结束标记。命令替换不是字符串转义工具，不拼接不可信文本形成命令。

\section{Python 参考程序}
\index{Python}\index{pow}
Python 的 int 为任意精度整数；\texttt{//} 是向下取整，负数时与 C++ 整数除法不同。\texttt{pow(a,e,m)} 做模幂；Python 3.8 起，\texttt{pow(a,-1,m)} 在可逆时求逆，不可逆会抛异常。\texttt{math.isqrt} 返回整数平方根，避免大整数浮点舍入。
\begin{lstlisting}[language=Python]
import math
import random
rng = random.Random(20260913)
a = rng.randrange(1, 100)
assert (-3) // 2 == -2
assert pow(3, -1, 7) == 5
assert math.isqrt(10**30) == 10**15
\end{lstlisting}
大量输入可用 \texttt{sys.stdin.buffer}，但一次 \texttt{read().split()} 会同时保存全部 token；大样例应估计 Python 对象开销。列表的 \texttt{b=a} 是同一个对象，\texttt{a[:]} 只复制一层；二维数组用列表推导构造，避免多行共享同一列表。

\chapter{语言基础}
\section{值、引用与调用边界}
\index{reference}
参数 \texttt{T} 按值，\texttt{const T\&} 只读借用，\texttt{T\&} 允许修改原对象。返回引用必须保证被引用对象仍存活，不能返回局部变量的引用。结构化绑定 \texttt{auto [x,y]} 默认复制，\texttt{auto\& [x,y]} 引用原元素；只读遍历常用 \texttt{const auto\&}。

\section{std::move：允许移动，不是立即移动}\label{infra-move}
\index{std::move}
\texttt{std::move(x)} 本身是转换表达式，使后续重载选择可以使用移动构造或移动赋值；它自己不搬运数据。对 const 对象使用 move 不会去掉 const，常见容器因此仍走复制。除有更强保证的特例外，标准库对象移动后有效但状态未指定：可销毁、重新赋值，或调用满足其前提的操作，不假定内容保持原样或必定为空。
\begin{lstlisting}
vector<int> a{1, 2, 3};
vector<int> b = std::move(a);
a.clear();
a.push_back(7);
const vector<int> c{4, 5};
vector<int> d = std::move(c);
\end{lstlisting}
上例 b 得到原序列，d 从 const c 复制。调用接受值参数的卷积函数且之后不再需要输入时，传 \texttt{std::move(a)} 可以避免额外副本；函数参数为 const 引用时不会因此取得所有权。返回本地变量通常直接 \texttt{return a;}；多加 move 可能失去具名返回值优化机会。依据：\href{https://eel.is/c++draft/forward}{标准草案 move 定义}与\href{https://eel.is/c++draft/lib.types.movedfrom}{移动后状态}。

\section{vector、array、deque 与字符串}
\index{vector}\index{array}\index{deque}\index{string}
vector 连续存储，随机访问常数时间，尾部追加摊还常数时间。\texttt{reserve} 只改容量，不增加元素；\texttt{resize} 才改变 size。重分配使原指针、引用和迭代器失效；erase 使被删位置及之后失效。不要保存 vector 节点引用后再 push 新节点；需要跨扩容引用时保存下标并重新访问。\href{https://eel.is/c++draft/vector.modifiers}{失效规则来源}。

array 的长度为编译期常量；局部基本类型元素需用 \texttt{\{\}} 初始化。deque 支持两端操作且可下标访问，但整体不连续；修改后不要盲目保留迭代器。vector<bool> 的元素可能是代理对象，不等同于普通 bool 引用。

string 适合字符序列；\texttt{substr(pos,len)} 第二个参数是长度。\texttt{find} 未找到返回 npos，不能直接转 int 再判断。\texttt{getline} 读整行，与 \texttt{cin>>} 混用时处理遗留换行。字符串反复头部删除可能线性搬移；需要复杂序列操作见分册 rope 与序列树。

\section{set、map 与哈希表}
\index{set}\index{multiset}\index{map}\index{unordered-map}
set/map 去重，multiset/multimap 允许等价键。比较器必须是严格弱序；不能用 less\_equal 假装多重集合，也不直接用不传递的 epsilon 比较做排序或树比较。map 的下标访问会插入默认值，只查询用 find。

multiset 的 \texttt{erase(key)} 会删掉全部等价元素；只删一个时先 find，再 erase 该迭代器。树容器的 lower\_bound 用成员函数，通用算法在非随机访问迭代器上可能线性走步。关联容器的键不可原地修改为破坏排序的值。\href{https://eel.is/c++draft/associative.reqmts}{关联容器要求}。

unordered\_map 的普通查改平均常数时间，最坏可线性；rehash 会使迭代器失效。需要容量控制时看其 reserve/max\_load\_factor 接口，不把这些名字照搬到 PBDS。GNU 的 gp/cc 哈希表、选型和范围见本册末的分册索引。

\section{queue、stack 与优先队列}
\index{queue}\index{stack}\index{priority-queue}
queue 先入先出，stack 后入先出；取 front/top 和 pop 前检查非空。标准 priority\_queue 默认大根堆，pop 不返回值，先复制 top 再 pop。它没有按句柄修改和合并接口。
\begin{lstlisting}
priority_queue<int, vector<int>, greater<int>> q;
q.push(7);
q.push(2);
int x = q.top();
q.pop();
\end{lstlisting}
需要修改已入堆的键、删除指定元素或合并时，使用数据结构册的 PBDS 配对堆；不要混用两种 priority\_queue 的接口或句柄保证。

\section{排序、查找与区间算法}
\index{sort}\index{lower-bound}\index{unique}\index{nth-element}
算法区间通常为半开区间。sort 要求严格弱序，stable\_sort 保留等价元素的原顺序；nth\_element 只保证第 k 位及两侧分区，平均线性，不等于完整排序。lower/upper\_bound 要求对应有序或分区前提，并使用一致比较规则。\href{https://eel.is/c++draft/alg.sorting}{算法要求}。
\begin{lstlisting}
sort(a.begin(), a.end());
a.erase(unique(a.begin(), a.end()), a.end());
auto it = lower_bound(a.begin(), a.end(), 7);
if (it != a.end()) cout << *it << '\n';
\end{lstlisting}
unique 只合并相邻等价值，返回逻辑终点；它不改容器 size。remove/remove\_if 同理，配合 erase；C++20 可用适用容器的 erase/erase\_if。reverse、rotate、merge、copy 的目标区间必须有效，不能把未扩容的 vector.begin() 当作足够大的输出区间。

\section{pair、tuple、optional 与回调}
\index{pair}\index{tuple}\index{optional}\index{function}
pair/tuple 默认按字典序比较，适合把值与唯一编号组合。optional 区分无解与合法零值，先判断再解引用。\texttt{std::function} 可保存不同可调用对象，但有类型擦除开销；泛型 lambda 或模板回调更容易内联。递归 lambda 捕获的引用不能活得比被引用对象更久。

\section{数值、位操作和随机数}
\index{numeric}\index{bitset}\index{random}
accumulate 的初值决定累加类型，64 位求和用 \texttt{0LL}。iota 适合连续编号。gcd/lcm 仍有可表示范围前提，最小负数的绝对值应先转宽或转无符号幅值。浮点整数精确范围和 long double 有效位数用 numeric\_limits 查询；几何数值条件另看几何册。

C++20 的 popcount/countr\_zero 等用于无符号类型；标准 countr\_zero(0) 有定义。

GCC 的 \texttt{\_\_builtin\_ctz(0)} 不能调用。\texttt{1ULL<<k} 仍须保证 k 小于位宽。bitset 的大小固定，批量位运算不是与位数无关的常数操作。

mt19937/mt19937\_64 用固定种子可复现对拍；uniform\_int\_distribution 的两端都包含。\texttt{rng()\%n} 一般有偏差，不把它当作严格均匀抽样。随机算法的种子和期望复杂度约定按对应模板说明执行。

\section{GNU 扩展与查询范围}
\index{PBDS}\index{rope}
PBDS 有序树、gp/cc 哈希表、配对堆以及 \texttt{\_\_gnu\_cxx::rope} 在所属知识点详细介绍，本册提供集中页码跳转。rope 不叫 std::rope；复制版本的内存成本和已知失败用法不因它是库容器而消失。
